The local \ac{sa} from Assignment~2 shows parameter influence at one specific point in the parameter space --- the calibrated optimum. This is useful, but limited: perturbing one parameter at a time misses any interactions between parameters, and the results say nothing about how the model behaves outside the immediate neighborhood of the calibrated values.

Global \ac{sa} addresses both of these limitations by varying all parameters simultaneously across a defined range and decomposing the resulting output variance. The Sobol method \autocite{Saltelli.2002} provides two indices for each parameter: the first-order index $S_1$ measures how much of the total output variance is due to that parameter alone, and the total-order index $S_T$ includes contributions from all interactions involving that parameter. A large gap between $S_T$ and $S_1$ indicates that a parameter acts mainly through interactions with others.

Two configuration choices are investigated here. First, the parameter range: the full bounds from Assignment~1 sample the entire admissible space, while a narrow range around the calibrated values gives a picture closer to realistic uncertainty. Second, the objective function: \ac{nse} weights high-flow performance, while logNSE emphasises low flows through the log transformation. Running both combinations shows how these choices affect which parameters appear sensitive and which do not.
